\scanpage{209}
% archive sheet number: 208
\deleted{Les \unsure{d}}
% margin note written slanted
\margin{\unsure{compatibilité} à vérifier\ldots}
\unsure{Modulo} identification des $d_2$ par morphismes canoniques (de nature ``motivique'') on trouve que la dimension des $E_\infty^{pq}$ est indépendante de $\unsure{\ell|l}$. De plus, les \deleted{valeurs propres} pol. car. de Frobenius \inserted{inverse} opérant sont à coeff. entiers entiers indépendants de $\ell$.

Utilisons la suite exacte
\[
\unsure{\bullet|0}\longrightarrow H^{i-1}(X_{\bar\eta})_I\longrightarrow H^i(\unsure{X_\eta|X_{\bar\eta}})\unsure{-|\longrightarrow}H^i(X_{\bar\eta})^I\unsure{-|\longrightarrow}0
\tag*{(\unsure{6|a})}
\]
Nous \deleted{\unsure{aurons}} \inserted{obtenons} alors des informations sur $H^{\unsure{*|\bullet}}(X_{\bar\eta})^I$ et $H^*(X_{\bar\eta})_I$, savoir

a) \unsure{leurs dimensions} indépendantes de $\ell$ $(\neq p)$

b) Les pol. car. de frobenius inverse opérant sur eux sont à coeff. entiers indépendants de $\ell$, et leurs val. propres \unsure{sont|ont} de val. absolues ${\unsure{q|1}}^{i/2}$ dans $H^i(X_{\bar\eta})^I$ et $H^i(X_{\bar\eta})_I$.

\deleted{Il suffit de le prouver pour $i\leq n$, car si c'est \unsure{vrai pour $i$, c'est} vrai pour $2n-i$ par dualité.}

On procède par récurrence sur $i$. Pour $i=0$, on a $H^0(X_\eta)\simeq H^0(X_{\bar\eta})^I$, on gagne. Pour $i>0$, \unsure{on a} \deleted{$\unsure{H^{i-1}(X_{\bar\eta})|H^i(X_{\bar\eta})}$} \unsure{utilise|utilisé} \unsure{(6)|(a)}, \deleted{\unsure{par}} et l'hyp. de récurrence.
